\documentclass[11pt,a4paper]{article}
% =====================================================================
%  Shared preamble for the Part V lecture notes (lectures 19-22).
%
%  These four lectures sit outside Geron, so the notes ARE the primary
%  source and are examined on the same terms as the textbook parts.
%  Everything here is chosen to make that readable on paper: one column,
%  generous measure, and the same six-colour palette as the slides so a
%  student moving between the two is not re-learning what red means.
%
%  Build with pdflatex (twice, for the toc and the refs):
%      cd notes && latexmk -pdf lecture-19.tex
%  or  make -C notes
% =====================================================================
\usepackage[T1]{fontenc}
\usepackage[utf8]{inputenc}
\usepackage{newtxtext,newtxmath}      % Times-like: prints well, reads well
\usepackage[scaled=0.86]{inconsolata} % code, at a size that sits beside Times
\usepackage{microtype}
% newtx defines \Bbbk, \openbox, \proof and \endproof; amssymb and amsthm
% then redefine them and abort. Freeing the four names before those two
% packages load is the standard fix, and it costs nothing here: we use
% amsthm only for \newtheorem.
\let\Bbbk\relax
\usepackage{amsmath,amssymb}
\let\openbox\relax
\let\proof\relax
\let\endproof\relax
\usepackage{amsthm}
\usepackage{booktabs}
\usepackage{enumitem}
\usepackage{xcolor}
\usepackage{tcolorbox}
\usepackage{titlesec}
\usepackage{graphicx}
\usepackage[hidelinks]{hyperref}
\tcbuselibrary{skins,breakable}

% ---- the course palette, from assets/css/custom.css -------------------
\definecolor{aimlprimary}{HTML}{0B3D62}   % deep blue  - structure
\definecolor{aimlaccent} {HTML}{C0392B}   % brick red  - failures
\definecolor{aimlsuccess}{HTML}{14663A}   % green      - repairs
\definecolor{aimlmath}   {HTML}{6C3483}   % purple     - the derivation
\definecolor{aimlmuted}  {HTML}{4B5563}
\definecolor{aimlrule}   {HTML}{B0BCC7}
\definecolor{aimlsoft}   {HTML}{F4F7F9}

\usepackage[a4paper,margin=32mm,top=30mm,bottom=30mm]{geometry}
\setlength{\parskip}{0.45\baselineskip}
\setlength{\parindent}{0pt}
\linespread{1.06}

% ---- headings --------------------------------------------------------
\titleformat{\section}{\Large\bfseries\color{aimlprimary}}{\thesection}{0.7em}{}
\titleformat{\subsection}{\large\bfseries\color{aimlprimary}}{\thesubsection}{0.7em}{}
\titleformat{\subsubsection}{\bfseries\color{aimlmuted}}{}{0em}{}
\titlespacing*{\section}{0pt}{1.6\baselineskip}{0.5\baselineskip}
\titlespacing*{\subsection}{0pt}{1.2\baselineskip}{0.35\baselineskip}

% ---- boxes -----------------------------------------------------------
% One box per role, and the role is the colour. A student who has seen the
% slides already knows what each one means before reading a word of it.
\newtcolorbox{keybox}[1][]{
  breakable, enhanced, colback=aimlsoft, colframe=aimlprimary,
  boxrule=0.4pt, left=8pt, right=8pt, top=6pt, bottom=6pt,
  fonttitle=\bfseries\small, coltitle=white, colbacktitle=aimlprimary,
  attach boxed title to top left={yshift=-2mm, xshift=4mm},
  boxed title style={boxrule=0pt, arc=1pt}, #1}

\newtcolorbox{failbox}[1][]{
  breakable, enhanced, colback=white, colframe=aimlaccent,
  boxrule=0.9pt, left=8pt, right=8pt, top=6pt, bottom=6pt,
  fonttitle=\bfseries\small, coltitle=white, colbacktitle=aimlaccent,
  attach boxed title to top left={yshift=-2mm, xshift=4mm},
  boxed title style={boxrule=0pt, arc=1pt}, #1}

\newtcolorbox{derivbox}[1][]{
  breakable, enhanced, colback=white, colframe=aimlmath,
  boxrule=0.6pt, left=8pt, right=8pt, top=6pt, bottom=6pt,
  fonttitle=\bfseries\small, coltitle=white, colbacktitle=aimlmath,
  attach boxed title to top left={yshift=-2mm, xshift=4mm},
  boxed title style={boxrule=0pt, arc=1pt}, #1}

% An exercise. Deliberately the same shape as the notebook's `try` field:
% one change, and what should happen to the output.
\newtcolorbox{trybox}[1][]{
  breakable, enhanced, colback=white, colframe=aimlsuccess,
  boxrule=0.5pt, left=8pt, right=8pt, top=6pt, bottom=6pt,
  fonttitle=\bfseries\small, coltitle=white, colbacktitle=aimlsuccess,
  attach boxed title to top left={yshift=-2mm, xshift=4mm},
  boxed title style={boxrule=0pt, arc=1pt}, #1}

% ---- small semantics -------------------------------------------------
\newcommand{\scope}[1]{\textcolor{aimlmuted}{\small #1}}
\newcommand{\term}[1]{\textbf{#1}}
\newcommand{\code}[1]{\texttt{\small #1}}
\newcommand{\lecture}[1]{Lecture~#1}

\newtheorem{definition}{Definition}[section]
\newtheorem{proposition}[definition]{Proposition}

% ---- title block -----------------------------------------------------
% \notestitle{19}{Information retrieval: the lexical foundation}{SciFact (BEIR)}
\newcommand{\notestitle}[3]{%
  \thispagestyle{empty}
  \begin{flushleft}
    {\color{aimlmuted}\small\sffamily
     APPLICATIONS OF MACHINE LEARNING \quad$\cdot$\quad
     BSc Mathematics of Artificial Intelligence}\\[2mm]
    {\color{aimlrule}\rule{\textwidth}{0.8pt}}\\[4mm]
    {\color{aimlmuted}\large Lecture #1 \quad$\cdot$\quad Part V \quad$\cdot$\quad
     Extended notes}\\[3mm]
    {\Huge\bfseries\color{aimlprimary}\baselineskip=1.15\baselineskip #2\par}\vspace{4mm}
    {\color{aimlmuted}\large Dataset: #3}\\[3mm]
    {\color{aimlrule}\rule{\textwidth}{0.8pt}}
  \end{flushleft}
  \vspace{2mm}
  \begin{keybox}[title={Why these notes exist}]
    Lectures 19--22 sit outside G\'eron, the course's primary text. For those
    four lectures \emph{these notes are the primary source}, and they are
    examined on exactly the same terms as the textbook parts. Everything in
    the slide deck for this lecture is here, with the arguments written out at
    length rather than compressed onto a slide; the numbers are the ones the
    accompanying notebook prints, and every one of them is a measurement on
    the corpus named above rather than a figure quoted from a paper.
  \end{keybox}
  \vspace{1mm}
  {\small\tableofcontents}
  \vspace{2mm}
  \noindent{\color{aimlrule}\rule{\textwidth}{0.4pt}}
  \vspace{1mm}
}


\begin{document}
\notestitlefor{7}{I}{Ensembles and random forests}{Forest CoverType}{Chapter 6}

\section{The result we are carrying forward}

Lecture 6 built a decision tree on this dataset and ended with a measurement
rather than a score. Refit the same depth-8 tree, with the same hyperparameters
and the same seed, on twenty different 90\% subsamples of the same training
rows: test accuracy moves by a standard deviation of 0.25 points, and
\emph{9.1\% of the individual predictions change}. Only 71.4\% of test patches
get the same answer from all twenty.

A stable metric, an unstable model. In the language of Lecture 5's
decomposition this is \emph{variance}: sensitivity of the fitted function to
the particular training sample. The whole of today follows from one question:
what does averaging do to variance?

\begin{center}
\small
\begin{tabular}{lrr}
\toprule
\textbf{Forest CoverType} & \textbf{Test accuracy} &
\textbf{Conditions per justification} \\
\midrule
Always ``Lodgepole Pine'' & 48.8\% & 0 \\
Depth-8 tree, minimum leaf 20 & 73.0\% & 7.97 \\
Unconstrained tree & 82.6\% & 17.88 \\
\bottomrule
\end{tabular}
\end{center}

\subsection{Where many different models come from}

CART is deterministic: same data, same tree. To get twenty different trees you
must give them twenty different things to look at.

\begin{center}
\small
\begin{tabular}{p{0.54\textwidth}l}
\toprule
\textbf{Source of difference} & \textbf{Name} \\
\midrule
Different training algorithms on the same data & a voting classifier \\
The same algorithm on different random subsets of \emph{rows}, drawn with
  replacement & bagging \\
\ldots drawn without replacement & pasting \\
\ldots plus random subsets of \emph{columns} at each split & random forest \\
\ldots plus random \emph{thresholds} & extra-trees \\
\bottomrule
\end{tabular}
\end{center}

The list is in increasing order of how hard it works to make two members unlike
each other. That ordering is not an accident.

Why this lecture follows the one on trees: averaging attacks variance, and
trees are the high-variance model we have just measured; a tree is cheap to
fit, so two hundred are affordable; and an unconstrained tree has \emph{low
bias} --- which matters, because averaging does nothing whatever to bias, so
the members had better not have any to spare. That third point is why the
ensembles here use unconstrained members. It looks reckless after a lecture
spent limiting depth, and by the end of the next section it will look obvious.

\section{Derivation: the variance of an average of correlated predictors}

Forget trees for ten minutes. You have $n$ predictors of the same quantity,
each with variance $\sigma^2$, and you average them:
$\bar f = \frac{1}{n}\sum_t f_t$. What is $\operatorname{Var}(\bar f)$?

The answer everyone gives first is $\sigma^2/n$. That step is valid only when
the $f_t$ are \emph{uncorrelated} --- the variance of a sum is the sum of the
variances \emph{plus twice the sum of the covariances}. Dropping the
covariances is an assumption, and it is exactly the assumption that fails for
an ensemble of trees grown on the same dataset.

\subsection{Two, then $n$}

With $\operatorname{Var}(f_1) = \operatorname{Var}(f_2) = \sigma^2$ and
$\operatorname{Corr}(f_1,f_2) = \rho$:
\[ \operatorname{Var}\Bigl(\tfrac{f_1+f_2}{2}\Bigr)
   = \tfrac14\bigl[\sigma^2 + \sigma^2 + 2\rho\sigma^2\bigr]
   = \frac{\sigma^2(1+\rho)}{2}. \]
At $\rho = 0$ this is the familiar halving; at $\rho = 1$ it is $\sigma^2$,
because the two predictors are the same predictor and averaging achieves
nothing. The general case has both limits in it.

\begin{derivbox}[title={$n$ identically distributed predictors, common pairwise $\rho$}]
Assume each has variance $\sigma^2$ and
$\operatorname{Cov}(f_s, f_t) = \rho\sigma^2$ for $s \ne t$. Then
\[ \operatorname{Var}\Bigl(\sum_{t=1}^{n} f_t\Bigr)
   = \sum_{s=1}^{n}\sum_{t=1}^{n}\operatorname{Cov}(f_s, f_t), \]
a double sum of $n^2$ terms. Split them by whether $s = t$: $n$ diagonal terms
each $\sigma^2$, and $n(n-1)$ off-diagonal terms each $\rho\sigma^2$. Divide by
$n^2$:
\[ \operatorname{Var}(\bar f)
   = \frac{1}{n^2}\bigl[n\sigma^2 + n(n-1)\rho\sigma^2\bigr]
   = \frac{\sigma^2}{n} + \frac{n-1}{n}\rho\sigma^2, \]
and writing $\frac{n-1}{n} = 1 - \frac1n$ and gathering the $1/n$ terms,
\[ \boxed{\;\operatorname{Var}(\bar f)
   = \rho\sigma^2 + \frac{(1-\rho)\sigma^2}{n}\;} \]
\end{derivbox}

\emph{Identically distributed, not independent} --- that is the whole point.
$n$ trees grown on bootstrap samples of one dataset have the same distribution
and are certainly not independent: they were handed the same dataset.

\subsection{Reading the two terms}

\begin{center}
\small
\begin{tabular}{lll}
\toprule
\textbf{Term} & \textbf{Depends on $n$?} & \textbf{What it is} \\
\midrule
$\rho\sigma^2$ & no & the part every predictor shares \\
$(1-\rho)\sigma^2/n$ & yes, like $1/n$ & the part each predictor has to itself \\
\bottomrule
\end{tabular}
\end{center}

\begin{keybox}[title={The sentence to remember}]
Averaging destroys the independent component of the variance and leaves the
correlated one completely untouched.

And note what it does \emph{not} say. It says nothing about bias: an average of
$n$ biased predictors is exactly as biased as one of them.
\end{keybox}

Three limits are worth memorising. At $\rho = 0$ the average has variance
$\sigma^2/n$, the classical result. At $\rho = 1$ it has variance $\sigma^2$
and averaging is a no-op. As $n \to \infty$ it tends to $\rho\sigma^2$, the
floor --- and because the second term falls like $1/n$, that floor is reached
quickly: going from 1 member to 10 removes 90\% of the independent part, while
going from 100 to 1{,}000 removes 0.9\% more.

So the way to improve an ensemble is \emph{not more members}. It is less
correlated members. And --- the point that gets forgotten --- anything which
lowers $\rho$ by making individual members worse is a trade, not a free lunch,
because $\sigma^2$ appears in \emph{both} terms and can rise while $\rho$
falls.

\begin{keybox}[title={One formula, three algorithms}]
Bagging randomises the rows. A random forest also randomises the columns.
Extra-trees also randomises the thresholds. All three exist for one purpose: to
lower $\rho$, because $\rho\sigma^2$ is the floor that $n$ cannot reach below.
\end{keybox}

\section{Building them, and measuring $\rho$}

\term{Bagging} draws $m$ rows with replacement, fits a tree, and repeats.
\term{Pasting} is the same without replacement, producing less diverse members.
A \term{random forest} adds a random subset of features at every node --- the
classification default \code{max\_features="sqrt"} gives
$\lfloor\sqrt{54}\rfloor = 7$ of our 54 columns, drawn afresh at each node, so
a tree with 200 nodes draws 200 different subsets. \term{Extra-trees} goes one
step further: having chosen the candidate features, draw one threshold at
random for each rather than searching for the best. Note that Scikit-Learn's
extra-trees default is \code{bootstrap=False}, so the random thresholds
\emph{replace} the bootstrap rather than joining it.

\subsection{A validation set you get for nothing}

Each bagged member never sees about 37\% of the rows --- its \term{out-of-bag}
instances. A given row is missed by one draw with probability $1 - 1/m$, so by
all $m$ draws with probability
\[ \Bigl(1 - \frac{1}{m}\Bigr)^{m} \longrightarrow e^{-1} \approx 0.368, \]
hence about 63\% distinct rows in and 37\% held out. Score each member on its
own out-of-bag rows and you have a generalisation estimate with no validation
split and no extra fit.

Read it carefully, though. It is an average over members, each scored on a
different 37\%, so it estimates the generalisation of one member's
\emph{procedure}, not of the ensemble you ship. It is available only where
members see bootstrap samples. And it does not replace the test set --- it
replaces the validation set.

\begin{keybox}[title={A leakage detector you get free}]
If the out-of-bag score comes out much \emph{better} than the test score,
something in the pipeline saw the out-of-bag rows anyway --- a scaler fitted
outside the ensemble, or a feature built from the whole frame.
\end{keybox}

\subsection{Measuring $\rho$ needs something we do not have}

$\rho$ is the correlation between two members over the randomness of the whole
procedure, which includes \emph{which training set you were handed}. Two trees
grown on one dataset are correlated through that dataset --- so a single
training set cannot show it. Conditional on the data, the members are
independent by construction and $\rho$ would measure as zero.

CoverType has 581{,}012 rows, so we can cut 20 \emph{disjoint} training sets of
20{,}000 rows each --- 400{,}000 rows, no row reused --- and still keep 12{,}000
aside as a shared test set. At each test patch, decompose the variance of a
single member across all 400 fitted trees:
\[ \sigma^2 = \underbrace{\tau^2}_{\text{between training sets}}
            + \underbrace{w^2}_{\text{within a training set}},
   \qquad \rho = \frac{\tau^2}{\tau^2 + w^2}. \]
$\tau^2$ is the part every member of a given ensemble inherits from its data,
and it \emph{is} the floor $\rho\sigma^2$ --- not a coincidence, but the same
quantity written twice.

\begin{failbox}[title={One trap in that code, worth ten minutes of your life}]
\begin{verbatim}
>>> rnd.classes_                    # the ensemble's labels
array([1, 2, 3, 4, 5, 6, 7])
>>> rnd.estimators_[0].classes_     # one member's labels
array([0., 1., 2., 3., 4., 5., 6.])
\end{verbatim}
Every ensemble re-encodes the labels as \emph{positions} before handing them to
its members, so a member's \code{predict} returns positions, not cover types.
Compare them with \code{y\_test} directly and every member scores about 8\% ---
which looks like a terrible model rather than a bug.
\end{failbox}

\subsection{Does the formula actually hold?}

\begin{center}
\small
\begin{tabular}{lrrrr}
\toprule
\textbf{Ensemble} & $\sigma^2$ & $\rho$ &
\textbf{Var of the average of 20} & \textbf{\ldots predicted} \\
\midrule
Bagging        & 0.1542 & 0.072 & 0.01814 & 0.01819 \\
Random forest  & 0.1675 & 0.043 & 0.01508 & 0.01523 \\
Extra-trees    & 0.1648 & 0.067 & 0.01859 & 0.01867 \\
Extra-trees, \code{bootstrap=True} & 0.1735 & 0.043 & 0.01554 & 0.01570 \\
\bottomrule
\end{tabular}
\end{center}

Measured and predicted agree in the third decimal place on all four. The
derivation is a description of this dataset. Agreement this good is partly
luck --- the formula assumes every pair has the same $\rho$, which is only
approximately true --- so do not expect three decimals every time.

\subsection{Does each variant lower $\rho$?}

\begin{center}
\small
\begin{tabular}{lllr}
\toprule
\textbf{Ensemble} & \textbf{Rows} & \textbf{Columns / thresholds} & $\rho$ \\
\midrule
Bagging       & bootstrap & all 54, best threshold & 0.072 \\
Extra-trees (default) & all & 7 of 54, random threshold & 0.067 \\
Random forest & bootstrap & 7 of 54, best threshold & 0.043 \\
Extra-trees, \code{bootstrap=True} & bootstrap & 7 of 54, random & 0.043 \\
\bottomrule
\end{tabular}
\end{center}

Every variant is below bagging, so the claim survives. But read the last two
rows: adding random thresholds on top of a random forest moves $\rho$ from
0.043 to 0.043. \emph{The mechanisms are not additive.}

Feature subsampling is doing the work, cutting $\rho$ by 40\%. Random
thresholds are a \emph{substitute} for the bootstrap rather than an addition to
it --- extra-trees without a bootstrap lands at 0.067, near bagging; with one,
exactly where the forest is. And there is a floor to the floor: nothing we
tried gets below 0.043, because all four ensembles ultimately saw the same
20{,}000 rows.

\begin{keybox}[title={The formula's real message}]
$\rho$ is not a property of the algorithm. It is the share of the variance that
comes from \emph{which data you were given}, and no amount of algorithmic
randomness reaches it.
\end{keybox}

\begin{center}
\small
\begin{tabular}{lrrrr}
\toprule
\textbf{Ensemble} & \textbf{One member} & \textbf{Average of 20} &
\textbf{Removed} & \textbf{Floor $\rho\sigma^2$} \\
\midrule
Bagging       & 0.15453 & 0.01814 & 88.3\% & 0.01103 \\
Random forest & 0.16734 & 0.01508 & 91.0\% & 0.00721 \\
Extra-trees   & 0.16439 & 0.01859 & 88.7\% & 0.01098 \\
Extra-trees, boot & 0.17248 & 0.01554 & 91.0\% & 0.00740 \\
\bottomrule
\end{tabular}
\end{center}

Twenty members remove about nine tenths of a single member's variance. Two
hundred would remove roughly one twentieth more, and then stop --- exactly what
$(1-\rho)\sigma^2/n$ says.

\section{And what did that do to the accuracy?}

\begin{center}
\small
\begin{tabular}{lrr}
\toprule
\textbf{200 members, 48{,}000 training patches} & \textbf{10 members} &
\textbf{200 members} \\
\midrule
Bagging       & 87.7\% & 89.8\% \\
Extra-trees   & 86.7\% & 88.9\% \\
Random forest & 85.9\% & 88.4\% \\
Extra-trees, \code{bootstrap=True} & 85.2\% & 88.0\% \\
\midrule
The legible depth-8 tree & --- & 73.0\% \\
Always ``Lodgepole Pine'' & --- & 48.8\% \\
\bottomrule
\end{tabular}
\end{center}

Every ensemble beats the constrained tree by at least fifteen points and the
unconstrained tree by at least five. The last 190 members buy about two.

Now put the two tables beside each other:

\begin{center}
\small
\begin{tabular}{lrrr}
\toprule
\textbf{Ensemble} & $\rho$ & \textbf{One member's accuracy} &
\textbf{Ensemble accuracy} \\
\midrule
Bagging       & 0.072 & 78.8\% & 89.8\% \\
Extra-trees   & 0.067 & 77.8\% & 88.9\% \\
Random forest & 0.043 & 73.6\% & 88.4\% \\
Extra-trees, boot & 0.043 & 72.8\% & 88.0\% \\
\bottomrule
\end{tabular}
\end{center}

\emph{The ensemble with the highest $\rho$ is the most accurate.}

\begin{failbox}[title={That is not a contradiction --- read the formula again}]
The formula is about \emph{variance}. It says nothing about bias, and accuracy
depends on both.

Every mechanism that lowers $\rho$ does so by \emph{handicapping the members}.
A tree allowed only 7 of 54 columns per node is a worse tree: 73.6\% against
bagging's 78.8\%. And $\sigma^2$ appears in both terms, so lowering $\rho$
while raising $\sigma^2$ can leave you worse off.

On this dataset --- 54 columns of which perhaps a dozen carry most of the
signal --- throwing away 47 at every node costs more than the decorrelation is
worth. On a dataset with hundreds of weakly informative columns it is usually
the other way round, which is why the default exists.
\end{failbox}

\begin{keybox}[title={What the derivation predicted, and what it declined to}]
It predicted the shape of the accuracy curve, the variance of a 20-member
average to three decimals on four different ensembles, and that more members
would stop helping and roughly when. It did \emph{not} tell us which ensemble
would be most accurate, and it never claimed to.

A theory that predicts three things and declines to predict the fourth is not a
weak theory. A theory that predicts all four without being asked is usually
being misread.
\end{keybox}

\section{The bill}

\begin{failbox}[title={An ensemble has no per-prediction justification}]
The 200 trees hold \textbf{1{,}497{,}919 leaves} between them. The
justification for one prediction is now 200 decision paths and a vote.

The agency asked for a model whose every prediction comes with a human-readable
justification. We have built one that is 15.4 points more accurate and cannot
supply one. \emph{There is no technical fix.} The tools below recover
something, and what they recover is genuinely less than what was lost. Saying
so is part of the job.
\end{failbox}

\begin{center}
\small
\begin{tabular}{ll}
\toprule
\textbf{The agency asked for} & \textbf{An ensemble can offer} \\
\midrule
Why this parcel & nothing comparable \\
Which measurements the decision rests on, overall & feature importance \\
How confident the model is & the vote share, better calibrated than one leaf \\
Reproducibility of the answer & much improved --- that was the whole repair \\
\bottomrule
\end{tabular}
\end{center}

Three of four is not nothing. The one we cannot supply is the one the agency
wrote down first.

On the vote share: a single tree's \code{predict\_proba} returns the class
proportions of one leaf, piecewise constant and often built on a handful of
rows. An ensemble averages 200 such vectors from 200 different leaves, so the
result varies smoothly and is better calibrated --- the one item on the
agency's list that averaging improves rather than destroys. But better
calibrated is not calibrated: a forest's vote share is systematically pulled
towards the middle, because unanimity is rare.

\subsection{Feature importance, and a control}

\code{feature\_importances\_} sums, over every node that split on a feature,
the weighted impurity reduction that split achieved --- the $\Delta I$ of
Lecture 6, accumulated. It runs, it is fast, and the top of the list is
entirely sensible: elevation first, then the distances. Nothing invites
suspicion.

\begin{keybox}[title={Before trusting any ranking}]
Ask what it would look like on a variable you know carries no information. If
you cannot answer, add one and find out.
\end{keybox}

Add a column of uniform random numbers and refit. It ranks \textbf{tenth of
fifty-five} --- above 45 of the 54 real measurements, with an impurity
importance of 0.0399 against 0.2277 for elevation. Put that list in front of
the agency and you have asserted that a column of noise matters more than 45
things the surveyors actually recorded.

\begin{failbox}[title={Two structural biases, not an accident}]
Impurity importance is measured on the \emph{training} data, and a split that
reduces impurity on the rows that chose it will do so whether or not the
feature is informative.

And it favours \emph{high-cardinality} features: a continuous column offers
thousands of candidate thresholds where a 0/1 soil-type column offers one, so
more chances to find a lucky split means more accumulated $\Delta I$. Our decoy
is continuous and 44 of the 54 real columns are binary.
\end{failbox}

The repair is to break the column and see what breaks: shuffle one column of
the \emph{held-out} set, re-score, and report a spread over repeats.

\begin{center}
\small
\begin{tabular}{lrr}
\toprule
\textbf{Measurement} & \code{random\_decoy} & \textbf{Elevation} \\
\midrule
Impurity importance (training data) & 0.0399 & 0.2277 \\
Rank among the 55 columns, by impurity & 10th & 1st \\
Permutation importance (held out) & $+0.00085 \pm 0.00195$ & far above it \\
\bottomrule
\end{tabular}
\end{center}

Permutation importance puts the decoy at less than half a standard deviation
from zero --- the correct answer, which the impurity measure never had any way
of reaching. And the forest with the decoy scores 87.6\% against 88.4\%
without: adding a noise column costs a little accuracy and a lot of
credibility.

\begin{center}
\small
\begin{tabular}{lp{0.36\textwidth}}
\toprule
\textbf{Question} & \textbf{Can importance answer it?} \\
\midrule
Which measurements should the survey keep collecting? & yes --- this is what it
                                                        is for \\
Is the model using a variable it legally must not? & yes, as a screen \\
Why was this parcel refused? & no \\
Would removing this column hurt? & no --- remove it, refit, and measure \\
\bottomrule
\end{tabular}
\end{center}

Row three is the agency's actual question, and it is the row importance cannot
fill. Do not present it as though it can.

\section{Boosting, and what practitioners reach for}

Everything so far trains members in parallel and averages them. \term{Boosting}
trains them in \emph{sequence}, each correcting its predecessor. The members
are not identically distributed and are certainly not exchangeable, so
\emph{today's formula does not apply to boosting}. Bagging reduces variance;
boosting reduces bias. Two consequences follow from the sequence: boosting
cannot be parallelised across members, and it is much easier to overfit with.

\term{AdaBoost} fits a weak learner --- typically a depth-1 stump ---
increases the weight of the instances it got wrong, and repeats, predicting by
a weighted vote. \term{Gradient boosting} keeps the sequence but changes the
correction: each new learner is fitted to the \emph{residual errors} of the
ensemble so far, with \code{learning\_rate} scaling each contribution.

\begin{failbox}[title={The intuition that does not carry across}]
Adding trees to a bagged ensemble can only help, because $(1-\rho)\sigma^2/n$
falls monotonically. Adding trees to a \emph{boosted} ensemble eventually
hurts. Carrying the first intuition into the second is one of the commonest
mistakes with these models.
\end{failbox}

\code{HistGradientBoostingClassifier} bins each feature into at most 255
buckets before growing anything, so a split search costs $O(\text{bins})$
rather than $O(m)$. It is the algorithm LightGBM and XGBoost implement, handles
missing values and categorical columns natively, and on tabular data is the
usual first choice. On this dataset, at a budget that fits a lecture, it lands
several points below the bagged forest --- usual is not always.

\term{Stacking} trains a model to do the aggregation, with the base learners'
predictions as the input features of a \emph{blender}. The blender must be
trained on predictions the base learners made for instances they did not see;
\code{StackingClassifier} does this with cross-validation internally, which is
the only reason it is safe to use. Build it by hand and the blender sees
in-sample predictions, learns to trust whichever base learner overfits
hardest, and the ensemble comes out worse than its best member --- with no
error and no warning.

\begin{center}
\small
\begin{tabular}{p{0.42\textwidth}p{0.46\textwidth}}
\toprule
\textbf{Situation} & \textbf{Reach for} \\
\midrule
A high-variance base learner and time to spare & bagging or a random forest \\
Many weakly informative columns & a random forest --- feature subsampling earns
                                  its keep \\
Training time is the binding constraint & extra-trees \\
You want the strongest tabular model & histogram gradient boosting --- usually,
                                       and measure it against a forest anyway \\
You must justify individual predictions & none of them --- go back to one
                                          tree \\
\bottomrule
\end{tabular}
\end{center}

The last row is what this pair of lectures exists to teach, and it is the row
people skip.

\section{Did the averaging fix what we diagnosed?}

The diagnosis was \emph{instability}, not inaccuracy --- so test the thing that
was diagnosed, not the thing that is easy to measure. On 20 disjoint training
sets, a random forest's variance of one member's correctness is 0.16734 and of
the average of 20 is 0.01508: 91.0\% of the variance removed, against a floor
of 0.00721 it cannot reach below. Nine tenths of a single member's variance,
gone --- and we knew before starting both that it would stop and roughly where.

\begin{keybox}[title={The standing constraint, extended by two clauses}]
Split before anything is fitted. All preprocessing inside a \code{Pipeline}
passed to cross-validation. Nothing derived from the test set in the training
path. Fixed seeds. Per-fold scores, not just the mean. \textbf{Any importance
or explanation is computed on held-out data and reported with a control. Any
per-instance claim is checked for stability under refitting.}

Both new clauses come from a measurement in this lecture rather than from a
principle: the decoy at rank 10 of 55, and 9.1\% of predictions moving while
the accuracy did not.
\end{keybox}

\section{The notebook}

\begin{trybox}[title={What to change}]
Set \code{max\_features=None} on the random forest, which turns it back into
bagging. Watch $\rho$ rise, each member get better, and the ensemble accuracy
rise with it --- the trade of the whole lecture, in one keyword.
\end{trybox}

\section{Exercises}

\begin{enumerate}[leftmargin=1.6em,itemsep=4pt]
  \item For $n$ predictors each of variance $\sigma^2$ and pairwise correlation
        $\rho$, the variance of their average is
        $\rho\sigma^2 + \frac{1-\rho}{n}\sigma^2$. State the limit as
        $n \to \infty$, and what it means for ensembles. \hfill \scope{[5]}
  \item Extra-trees are usually faster to fit than a random forest and often no
        less accurate. Give the mechanism, in terms of the formula above.
        \hfill \scope{[4]}
  \item You set \code{bootstrap=False} and \code{oob\_score=True}. State what
        happens and why. \hfill \scope{[3]}
  \item Impurity-based feature importance ranks a random decoy column above
        several real ones. Explain how that can happen, and name the measure
        you would use instead. \hfill \scope{[5]}
  \item A forest improves accuracy over one tree by far less than the variance
        reduction would suggest. Is this a contradiction? Answer with a reason.
        \hfill \scope{[4]}
\end{enumerate}

\section*{Summary}
\addcontentsline{toc}{section}{Summary}

From the double sum of covariances,
$\operatorname{Var}(\bar f) = \rho\sigma^2 + (1-\rho)\sigma^2/n$: averaging
destroys the independent component of the variance and cannot touch the
correlated one. Read as a policy, that says more members buy less and less, and
the way to improve an ensemble is less correlated members.

Measured on 20 disjoint training sets, the formula predicted the variance of a
20-member average to three decimals on four different ensembles. $\rho$ fell
from 0.072 to 0.043 --- and the mechanisms turned out not to be additive:
random thresholds substitute for the bootstrap rather than adding to it, and
nothing reached below 0.043 because all four ensembles saw the same 20{,}000
rows. $\rho$ is not a property of the algorithm but of which data you were
given.

The lowest-$\rho$ ensemble was not the most accurate, because every mechanism
that lowers $\rho$ also weakens the members and $\sigma^2$ sits in both terms.
And the 15.4 points of accuracy cost 1{,}497{,}919 leaves and the
per-prediction justification the previous lecture was built around. Impurity
importance ranked a column of pure noise 10th of 55; a permutation measurement
on held-out data put it half a standard deviation from zero, where it belongs.

\end{document}
